\section{AIME II 2008}
        \begin{enumerate}
        	\item (\textit{AIME II 2008}) Let $N = 100^2 + 99^2 - 98^2 - 97^2 + 96^2 + \cdots + 4^2 + 3^2 - 2^2 - 1^2$, where the additions and subtractions alternate in pairs. Find the remainder when $N$ is divided by $1000$.


 
 \textbf{Answer: }100



\textbf{Solution 1}

Rewriting this sequence with more terms, we have

 
 \begin{align*} N &= 100^2 + 99^2 - 98^2 - 97^2 + 96^2 + 95^2 - 94^2 - 93^2 + 92^2 + 91^2 + \ldots - 10^2 - 9^2 + 8^2 + 7^2 - 6^2 - 5^2 + 4^2 + 3^2 - 2^2 - 1^2 \mbox{, and reordering, we get}\\ N &= (100^2 - 98^2) + (99^2 - 97^2) + (96^2 - 94^2)  + (95^2 - 93^2) + (92^2 - 90^2) + \ldots +  (8^2 - 6^2) + (7^2 - 5^2) +(4^2 - 2^2) + (3^2 - 1^2) \mbox{.} \end{align*}
Factoring this expression yields

 
 \begin{align*} N &= (100 - 98)(100 + 98) + (99 - 97)(99 + 97) + (96 - 94)(96 + 94) + (95 - 93)(95 + 93) + (92 - 90)(90 + 92) + \ldots + (8 - 6)(8 + 6) + (7 - 5)(7 + 5) + (4 - 2)(4 + 2) + (3 - 1)(3 + 1) \mbox{, leading to}\\ N &= 2(100 + 98) + 2(99 + 97) + 2(96 + 94) + 2(95 + 93) + 2(92 + 90) + \ldots + 2(8 + 6) + 2(7 + 5) + 2(4 + 2) + 2(3 + 1) \mbox{.} \end{align*}
Next, we get

 
 \begin{align*} N &= 2(100 + 98 + 99 + 97 + 96 + 94 + 95 + 93 + 92 + 90 + \ldots + 8 + 6 + 7 + 5 + 4 + 2 + 3 + 1 \mbox{, and rearranging terms yields}\\ N &= 2(100 + 99 + 98 + 97 + 96 + \ldots + 5 + 4 + 3 + 2 + 1) \mbox{.} \end{align*}
Then,

 
 \begin{align*} N &= 2\left(\frac{(100)(101)}{2}\right) \mbox{, and simplifying, we get}\\ N &= (100)(101) \mbox{, so}\\ N &= 10100 \mbox{.} \end{align*}
Dividing $10100$ by $1000$ yields a remainder of $\boxed{100}$.

 





\textbf{Solution 2}

Since we want the remainder when $N$ is divided by $1000$, we may ignore the $100^2$ term. Then, applying the \href{https://artofproblemsolving.com/wiki/index.php?title=Difference\_of\_squares}{difference of squares} factorization to consecutive terms,

 
 \begin{align*} N &= (99-98)(99+98) - (97-96)(97+96) + (95-94)(95 + 94) + \cdots + (3-2)(3+2) - 1 \\ &= \underbrace{197 - 193}_4 + \underbrace{189 - 185}_4 + \cdots + \underbrace{5 - 1}_4 \\  &= 4 \cdot \left(\frac{197-5}{8}+1\right) = \boxed{100} \end{align*}






\textbf{Solution 3}

By observation, we realize that the sequence 
 \[(a+3)^2 + (a+2)^2 - (a+1)^2 - a^2\] alternates every 4 terms. Simplifying, we get 
 \[(a+3)^2 + (a+2)^2 - (a+1)^2 - a^2 = 8a + 12\], turning $N$ into a arithmetic sequence with 25 terms, them being $1, 5, 9, \dots ,97$, as the series $8a + 12$ alternates every 4 terms.

 Applying the sum of arithmetic sequence formula, we get

 
 \begin{align*} N &= \frac{25\cdot((8\cdot1 + 12) + (8\cdot97 + 12))}{2} \\ &= \frac{25\cdot(20 + 788)}{2} = 10100 \end{align*}
So the answer would be 
 \[\frac{10100}{1000} = \boxed{100}\].

 
- erdaifuu

 





\textbf{Solution 4}

We can remove the $100^2$ since $100^2 \equiv 0 \pmod{1000}$ and use difference of squares to factor out the rest. This gives 
 \[(1)(99+98) + (-1)(96+97) + ... +(1)(3+2) +(-1)(1+0)\] Writing this another way, we get 
 \[(99(2) - 1) - (97(2)- 1) + (95(2) - 1) -  ... +(3(2)-1) - (1(2) -1)\] We know that the last one is negative because all the numbers before multiplying that are in the form $4a - 1$(eg. $99 = 25(4) - 1$) are positive. 

 Let $x = 99$. This makes the expression 
 \[(2x-1) - (2(x-2) - 1) + (2(x-4) - 1) - ... + (2(x-96) - 1) -   (2(x-98) - 1)\] This simplifies to 
 \[(2x-1) - (2x-5) + (2x-9) - ... +(2x - 193) - (2x-197)\] Since the first one is positive and the last one is negative, that means there are an even number of terms and using the associative property and the distributive property, all of the $2x$ terms cancel out. A consequence of this is that all of the positive integers turn negative and all the negative ones turn positive(eg. $-(2x-5) = -2x + 5$). 

 We are left with the sequence 
 \[-1+5-9+...-193+197\] We can notice the property that the number of terms in the sequence to a positive number n is equal to $(n+3)/4$, as well as the fact that every pair sums up to 4. Therefore the total number of terms is $(197+3) / 4 = 50$. Therefore, there are 25 pairs each summing up to 4, leaving us with $25(4) = \boxed{100}$.

 





\textbf{Solution 5}

We simply take the outer pairs and from there, we use the inside terms. That is, $N = 100^2 + 99^2 - 98^2 - 97^2 + 96^2 + \cdots + 4^2 + 3^2 - 2^2 - 1^2$becomes $N = (100^2 - 1^2) + (99^2 - 2^2) - (98^2 - 3^2) - (97^2 - 4^2) + (96^2 - 5^2) ...$and thus is reduced to$N = 101(99 + 97 - 95 - 93 + 81 + 89 - ... - 7 - 5 + 3 + 1)$since the entire sequence has 50 terms. Therefore, the last two terms must be positive as the above inner sequence repeats as follows : Positive - positive - negative - negative. It follows that $N = 101[8(12) + 4] = 101[100] = 10100$. So the requested answer is $\boxed{100}$. ~elpianista227

 Note: Alternatively, we group the 50 odd numbers into 4's to obtain 12.5 groups, each have a consequent sum of 8, so 100 total. This works due to the values being the sum of two arithmetic series. 

 




	\item (\textit{AIME II 2008}) Rudolph bikes at a constant rate and stops for a five-minute break at the end of every mile. Jennifer bikes at a constant rate which is three-quarters the rate that Rudolph bikes, but Jennifer takes a five-minute break at the end of every two miles. Jennifer and Rudolph begin biking at the same time and arrive at the $50$-mile mark at exactly the same time. How many minutes has it taken them?


 
 \textbf{Answer: }620



\textbf{Solution 1}

Let Rudolf bike at a rate $r$, so Jennifer bikes at the rate $\dfrac 34r$. Let the time both take be $t$. 

 Then Rudolf stops $49$ times (because the rest after he reaches the finish does not count), losing a total of $49 \cdot 5 = 245$ minutes, while Jennifer stops $24$ times, losing a total of $24 \cdot 5 = 120$ minutes. The time Rudolf and Jennifer actually take biking is then $t - 245,\, t-120$ respectively. 

 Using the formula $r = \frac dt$, since both Jennifer and Rudolf bike $50$ miles,

 
 \begin{align}r &= \frac{50}{t-245}\\ \frac{3}{4}r &= \frac{50}{t-120} \end{align}Substituting equation $(1)$ into equation $(2)$ and simplifying, we find 

 
 \begin{align*}50 \cdot \frac{3}{4(t-245)} &= 50 \cdot \frac{1}{t-120}\\ \frac{1}{3}t &= \frac{245 \cdot 4}{3} - 120\\ t &= \boxed{620}\ \text{minutes} \end{align*}






\textbf{Solution 2}

Let the total time that Jennifer and Rudolph bike, including rests, to be $t$ minutes.  Furthermore, let Rudolph's biking rate be $r$ so Jennifer's biking rate is $\frac{3}{4}r$. Note that Rudolf takes 49 breaks, taking $49\cdot 5$ minutes, and Jennifer takes 24 breaks, taking $24\cdot 5$ minutes.  Since they both reach the 50 mile mark, then by $d=rt$, the rate times time taken for Rudolph and Jennifer must equal.  Hence, we disregard the breaks from the total time taken and get the equation  
 \[r(t-49\cdot 5)=\frac{3}{4}r(t-24\cdot 5),\]yielding $t=\boxed{620}$.

 





\textbf{Solution 3 (Similar to Solution 1, 1 and 2 combined)}

Recall $d=rt$. So, $d = r(t-(49\cdot 5))$ where $r$ is the speed of Rudolph, and $d=\frac{3}{4}r(t-(24 \cdot 5))$, describing Jennifer analogously. We amalgamate the two equations and reduce $r$ to obtain $t-245 = \frac{3}{4}t - 90$, so $\frac{1}{4}t = 155$, and $t = \boxed{620}$. Here, $r$ is unimportant, saving time, just like the simple solution stated previously.

 




	\item (\textit{AIME II 2008}) A block of cheese in the shape of a rectangular solid measures $10$ cm by $13$ cm by $14$ cm. Ten slices are cut from the cheese. Each slice has a width of $1$ cm and is cut parallel to one face of the cheese. The individual slices are not necessarily parallel to each other. What is the maximum possible volume in cubic cm of the remaining block of cheese after ten slices have been cut off?


 
 \textbf{Answer: }729



\textbf{Solution}

Let the lengths of the three sides of the rectangular solid after the cutting be $a,b,c$, so that the desired volume is $abc$. Note that each cut reduces one of the dimensions by one, so that after ten cuts, $a+b+c = 10 + 13 + 14 - 10 = 27$. By \href{https://artofproblemsolving.com/wiki/index.php?title=AM-GM}{AM-GM}, $\frac{a+b+c}{3} = 9 \ge \sqrt[3]{abc} \Longrightarrow abc \le \boxed{729}$. Equality is achieved when $a=b=c=9$, which is possible if we make one slice perpendicular to the $10$ cm edge, four slices perpendicular to the $13$ cm edge, and five slices perpendicular to the $14$ cm edge.

 





\textbf{Solution 2}

A more intuitive way to solve it is by seeing that to keep the volume of the rectangular cheese the greatest, we must slice the cheese off to decrease the greatest length of the cheese (this is easy to check). Here are the ten slices: 

 ${10, 13, 14} \rightarrow {10, 13, 13} \rightarrow {10, 12, 13} \rightarrow {10, 12, 12} \rightarrow {10, 11, 12} \rightarrow {10, 11, 11} \rightarrow {10, 10, 11} \rightarrow {10, 10, 10} \rightarrow {9, 10, 10} \rightarrow {9, 9, 10} \rightarrow {9, 9, 9}.$

 So the greatest possible volume is thus $9 \times 9 \times 9 = \boxed{729}$

 




	\item (\textit{AIME II 2008}) There exist $r$ unique nonnegative integers $n_1 > n_2 > \cdots > n_r$ and $r$ integers $a_k$ ($1\le k\le r$) with each $a_k$ either $1$ or $- 1$ such that

 \[a_13^{n_1} + a_23^{n_2} + \cdots + a_r3^{n_r} = 2008.\]
Find $n_1 + n_2 + \cdots + n_r$.


 
 \textbf{Answer: }021



\textbf{Solution}

In base $3$, we find that $\overline{2008}_{10} = \overline{2202101}_{3}$. In other words,

 $2008 = 2 \cdot 3^{6} + 2 \cdot 3^{5} + 2 \cdot 3^3 + 1 \cdot 3^2 + 1 \cdot 3^0$In order to rewrite as a sum of perfect powers of $3$, we can use the fact that $2 \cdot 3^k = 3^{k+1} - 3^k$:

 $2008 = (3^7 - 3^6) + (3^6-3^5) + (3^4 - 3^3) + 3^2 + 3^0 = 3^7 - 3^5 + 3^4 - 3^3 + 3^2 + 3^0$The answer is $7+5+4+3+2+0 = \boxed{021}$.

 Note : Solution by bounding is also possible, namely using the fact that $1+3+3^2 + \cdots + 3^{n} =\frac{3^{n+1}-1}{2}.$

 





\textbf{Solution 2}

Notice $2008$ has $7$ digits in base 3. Therefore, we can add $1093=1111111_3$ to get $3101=11020212_3$. Subtracting $1093$ back from this, we find $2008 = 1 0 (-1) 1 (-1) 1 0 1_3$. We need to sum the exponents of the non-zero digits of this number, so our answer is $7+5+4+3+2+0=\boxed{021}$

 




	\item (\textit{AIME II 2008}) In trapezoid $ABCD$ with $\overline{BC}\parallel\overline{AD}$, let $BC = 1000$ and $AD = 2008$. Let $\angle A = 37^\circ$, $\angle D = 53^\circ$, and $M$ and $N$ be the midpoints of $\overline{BC}$ and $\overline{AD}$, respectively. Find the length $MN$.


 
 \textbf{Answer: }504


	\item (\textit{AIME II 2008}) The sequence $\{a_n\}$ is defined by

 \[a_0 = 1,a_1 = 1, \text{ and } a_n = a_{n - 1} + \frac {a_{n - 1}^2}{a_{n - 2}}\text{ for }n\ge2.\]
The sequence $\{b_n\}$ is defined by

 \[b_0 = 1,b_1 = 3, \text{ and } b_n = b_{n - 1} + \frac {b_{n - 1}^2}{b_{n - 2}}\text{ for }n\ge2.\]
Find $\frac {b_{32}}{a_{32}}$.


 
 \textbf{Answer: }561



\textbf{Solution}

Rearranging the definitions, we have
 \[\frac{a_n}{a_{n-1}} = \frac{a_{n-1}}{a_{n-2}} + 1,\quad \frac{b_n}{b_{n-1}} = \frac{b_{n-1}}{b_{n-2}} + 1\]from which it follows that $\frac{a_n}{a_{n-1}} = 1+ \frac{a_{n-1}}{a_{n-2}} = \cdots = (n-1) + \frac{a_{1}}{a_0} = n$ and $\frac{b_n}{b_{n-1}} = (n-1) + \frac{b_{1}}{b_0} = n+2$. These \href{https://artofproblemsolving.com/wiki/index.php?title=Recursion}{recursions}, $a_{n} = na_{n-1}$ and $b_{n} = (n+2)b_{n-1}$, respectively, correspond to the explicit functions $a_n = n!$ and $b_n = \frac{(n+2)!}{2}$ (after applying our initial conditions). It follows that $\frac{b_{32}}{a_{32}} = \frac{\frac{34!}{2}}{32!} = \frac{34 \cdot 33}{2} = \boxed{561}$.

 From this, we can determine that the sequence $\frac {b_n}{a_n}$ corresponds to the \href{https://artofproblemsolving.com/wiki/index.php?title=Triangular\_number}{triangular numbers}.

 




	\item (\textit{AIME II 2008}) Let $r$, $s$, and $t$ be the three roots of the equation

 \[8x^3 + 1001x + 2008 = 0.\]
Find $(r + s)^3 + (s + t)^3 + (t + r)^3$.


 
 \textbf{Answer: }753



\textbf{Solution 1}

By \href{https://artofproblemsolving.com/wiki/index.php?title=Vieta\%27s\_formulas}{Vieta's formulas}, we have $r + s + t = 0$ so $t = -r - s.$ Substituting this into our problem statement, our desired quantity is 
 \[(r + s)^3 - r^3 - s^3 = 3r^2s + 3rs^2 = 3rs(r + s).\] Also by \href{https://artofproblemsolving.com/wiki/index.php?title=Vieta\%27s\_formulas}{Vieta's formulas} we have 
 \[rst = -rs(r + s) = -\dfrac{2008}{8} = -251\] so negating both sides and multiplying through by 3 gives our answer of $\boxed{753}.$

 





\textbf{Solution 2}

By \href{https://artofproblemsolving.com/wiki/index.php?title=Vieta\%27s\_formulas}{Vieta's formulas}, we have $r+s+t = 0$, and so the desired answer is $(r+s)^3 + (s+t)^3 + (t+r)^3 = (0-t)^3 + (0-r)^3 + (0-s)^3 = -(r^3 + s^3 + t^3)$. Additionally, using the factorization
 \[r^3 + s^3 + t^3 - 3rst = (r+s+t)(r^2 + s^2 + t^2 - rs - st - tr) = 0\]we have that $r^3 + s^3 + t^3 = 3rst$. By Vieta's again, $rst = \frac{-2008}8 = -251 \Longrightarrow -(r^3 + s^3 + t^3) = -3rst = \boxed{753}.$

 





\textbf{Solution 3}

Vieta's formulas gives $r + s + t = 0$. Since $r$ is a root of the polynomial, $8r^3 + 1001r + 2008 = 0\Longleftrightarrow - 8r^3 = 1001r + 2008$, and the same can be done with $s,\ t$. Therefore, we have
 \begin{align*}8\{(r + s)^3 + (s + t)^3 + (t + r)^3\} &= - 8(r^3 + s^3 + t^3)\\ &= 1001(r + s + t) + 2008\cdot 3 = 3\cdot 2008\end{align*}
yielding the answer $\boxed{753}$.

 Also, Newton's Sums yields an answer through the application.\href{https://artofproblemsolving.comhttp://www.artofproblemsolving.com/Wiki/index.php/Newton's\_Sums}{http://www.artofproblemsolving.com/Wiki/index.php/Newton's\_Sums}

 





\textbf{Solution 4}

Expanding, you get:
 \[r^3 + 3r^2s + 3s^2r +s^3 +\] 
 \[s^3 + 3s^2t + 3t^2s +t^3 +\]
 \[r^3 + 3r^2t + 3t^2r +t^3\]
 \[= 2r^3 + 2s^3 + 2t^3 + 3r^2s + 3s^2r + 3s^2t + 3t^2s + 3r^2t + 3t^2r\]This looks similar to $(r+s+t)^3 = r^3 + s^3 + t^3 + 3r^2s + 3s^2r + 3s^2t + 3t^2s + 3r^2t + 3t^2r + 6rst$Substituting:
 \[(r+s+t)^3 - 6rst + r^3+s^3+t^3 = (r + s)^3 + (s + t)^3 + (t + r)^3\]Since $r+s+t = 0$, 
 \[(r+s)^3 + (s+t)^3 + (t+r)^3 = (0-t)^3 + (0-r)^3 + (0-s)^3 = -(r^3 + s^3 + t^3)\]Substituting, we get 
 \[(r+s+t)^3 - 6rst + r^3+s^3+t^3 =  -(r^3 + s^3 + t^3)\]or, 
 \[0^3 - 6rst + r^3+s^3+t^3 =  -(r^3 + s^3 + t^3) \implies 2(r^3 + s^3 + t^3) = 6rst\]We are trying to find $-(r^3 + s^3 + t^3)$.Substituting:
 \[-(r^3 + s^3 + t^3) = -3srt = \frac{-2008*3}{8} = \boxed{753}.\]

 





\textbf{Solution 5}

Write $(r+s)^3+(s+t)^3+(t+r)^3=-(r^3+s^3+t^3)$ and let $f(x)=8x^3+1001x+2008$. Then 
 \[f(r)+f(s)+f(t)=8(r^3+s^3+t^3)+1001(r+s+t)+6024=8(r^3+s^3+t^3)+6024=0.\] Solving for $r^3+s^3+t^3$ and negating the result yields the answer $\boxed{753}.$

 





\textbf{Solution 6}

Here by \href{https://artofproblemsolving.com/wiki/index.php?title=Vieta\%27s\_formulas}{Vieta's formulas}:$r+s+t = 0$ --(1)

 $rst = \frac{-2008}{8} = -251$ --(2)

 By the factorisation formula:Let $a = r+s$, $b = s+t$, $c = t+r$, $a^3+b^3+c^3-3abc = (a+b+c)(a^2+b^2+c^2-ab-bc-ca) = 0$ (By (1))

 So 
 \[a^3+b^3+c^3 = 3abc = 3(r+s)(s+t)(t+r) = 3(-t)(-r)(-s) = 3[-(-251)] = \boxed{753}.\]

 





\textbf{Solution 7}

Let's construct a polynomial with the roots $(r+s), (s+t),$ and $(t+r)$.

 sum of the roots:

 $=2(r+s+t)=2\cdot0=0$

 pairwise product of the roots:

 $(r+s)(s+t)+(s+t)(t+r)+(t+r)(r+s)=r^2+s^2+t^2+3(rs+st+tr)$

 $=(r+s+t)^2+rs+st+tr=0+\frac{1001}{8}$

 product of the roots:

 $(r+s)(s+t)(t+r)=r^2t+r^2s+s^2r+s^2t+t^2r+t^2s+3rst$

 $=(rs+st+tr)(r+s+t)-3rst+2rst=-rst=-\frac{2008}{8}$

 thus, the polynomial we get is

 $x^3+\frac{1001}{8}x+-\frac{2008}{8}=0$

 as $(r+s), (s+t),$ and $(t+r)$ are roots of this polynomial, we know that (using power reduction)

 $(r+s)^3+\frac{1001}{8}(r+s)-\frac{2008}{8}=0$

 $(s+t)^3+\frac{1001}{8}(s+t)-\frac{2008}{8}=0$

 $(t+r)^3+\frac{1001}{8}(t+r)-\frac{2008}{8}=0$

 adding all of the equations up, we see that

 $(r+s)^3+(s+t)^3+(t+r)^3=3\cdot\frac{2008}{8}-\frac{1001}{8}(2r+2s+2t)=251(3)+0=\boxed{753}$

 





\textbf{Solution 8}

We want to find what is $-(r^3+s^3+t^3)$ which reminds us of Newton sum. So we can see that $8S_3+0\cdot S_2+1001\cdot S_1+3\cdot 2008=0$ Notice that $S_1=0$ so it is just $S_3=-\frac{2008\cdot 3}{8}=-753$, the desired answer is $\boxed{753}$

 





\textbf{Solution 9}

This solution uses Vietas, as with everyone else's solution.Expanding the expression we get

 
 \[(r+s)^3+(s+t)^3+(t+r)^3 = r^3+3r^2s+3rs^2+\dots +3s^2t+3ts^2+t^3\]

 Seeing the cubes, we try to find a $(r+s+t)^3$ and upon doing so, we get

 
 \[(r+s)^3+(s+t)^3+(t+r)^3=(r+s+t)^3-6rst+(r^3+s^3+t^3)\]

 Recall that $a^3+b^3+c^3-3abc=(a+b+c)(a^2+b^2+c^2-ab-bc-ca)$. Thus, we get

 
 \[(r+s)^3+(s+t)^3+(t+r)^3=(r+s+t)^3-3rst+(r+s+t)(r^2+s^2+t^2-rs-st-tr)\]

 Plugging in $(r+s+t)=0$ we get

 
 \[(r+s)^3+(s+t)^3+(t+r)^3=0-3rst+0=-3\cdot -251=\boxed{753}\]

 





\textbf{Solution 10}

$8x^3+1001x+2008=0$

 We want to find $(r+s)^3+(s+t)^3+(t+r)^3.$ Let's call this result n.

 From vieta's formulas, we find that $r+s+t=-0/8=0$, $rs+st+tr=1001/8$, and $rst=-2008/8=-251.$

 Expanding and rearranging gives us $n=(r+s)^3+(s+t)^3+(t+r)^3=r^3+3r^2s+3rs^2+s^3+s^3+3s^2t+3st^2+t^3+t^3+3t^2r+3tr^2+r^3=2r^3+2s^3+3r^2s+3rs^2+3s^2t+3st^2+3t^2r+3tr^2=3(r^3+s^3+t^3-(r^2s+rs^2+s^2t+st^2+t^2r+tr^2))-(r^3+s^3+t^3)=3((r+s+t)(r^2+s^2+t^2))-((r+s+t)^3-3(r^2s+rs^2+s^2t+st^2+t^2r+tr^2)-6rst)=3(r+s+t)((r+s+t)^2-2(rs+st+tr))-((r+s+t)^3-3(r^2s+rs^2+s^2t+st^2+t^2r+tr^2)-6rst)$

 Let $k=r^2s+rs^2+s^2t+st^2+t^2r+tr^2$

 $k=r^2s+rs^2+s^2t+st^2+t^2r+tr^2=(r+s+t)(r^2+s^2+t^2)-(r^3+s^3+t^3)=(r+s+t)((r+s+t)^2-2(rs+st+tr))-((r+s+t)^3-3k-6rst)=(0)(0^2-2(1001/8))-(0^3-3k-6(-251))=0-(0-3k+1506)=3k-1506$

 Solving gives us $k=753$

 $n=3(r+s+t)((r+s+t)^2-2(rs+st+tr))-((r+s+t)^3-3(r^2s+rs^2+s^2t+st^2+t^2r+tr^2)-6rst)=3(0)(0^2-2(1001/8))-(0^3-3k-6(-251))=0-(0-3(753)+1506)=753$

 Therefore, the answer is $753.$

 Also note that this is the only solution that still would have worked effectively if $r+s+t$ was nonzero.

 




	\item (\textit{AIME II 2008}) Let $a = \pi/2008$. Find the smallest positive integer $n$ such that

 \[2[\cos(a)\sin(a) + \cos(4a)\sin(2a) + \cos(9a)\sin(3a) + \cdots + \cos(n^2a)\sin(na)]\]
is an integer.


 
 \textbf{Answer: }251



\textbf{Solution 1}

By the \href{https://artofproblemsolving.com/wiki/index.php?title=Trigonometric\_identity}{product-to-sum identities}, we have that $2\cos a \sin b = \sin (a+b) - \sin (a-b)$. Therefore, this reduces to a telescope series: 
 \begin{align*} \sum_{k=1}^{n} 2\cos(k^2a)\sin(ka) &= \sum_{k=1}^{n} [\sin(k(k+1)a) - \sin((k-1)ka)]\\ &= -\sin(0) + \sin(2a)- \sin(2a) + \sin(6a) - \cdots - \sin((n-1)na) + \sin(n(n+1)a)\\ &= -\sin(0) + \sin(n(n+1)a) = \sin(n(n+1)a) \end{align*}


 Thus, we need $\sin \left(\frac{n(n+1)\pi}{2008}\right)$ to be an integer; this can be only $\{-1,0,1\}$, which occur when $2 \cdot \frac{n(n+1)}{2008}$ is an integer. Thus $1004 = 2^2 \cdot 251 | n(n+1) \Longrightarrow 251 | n, n+1$. We know that $n$ cannot be $250$ as $250$ isn't divisible by $4$, so 1004 doesn't divide $n(n+1) = 250 \cdot 251$. Therefore, it is clear that $n = \boxed{251}$ is the smallest such integer.

 





\textbf{Solution 2}

We proceed with complex trigonometry. We know that for all $\theta$, we have $\cos \theta = \dfrac{1}{2} \left( z + \dfrac{1}{z} \right)$ and $\sin \theta = \dfrac{1}{2i} \left( z - \dfrac{1}{z} \right)$ for some complex number $z$ on the unit circle. Similarly, we have $\cos n \theta = \dfrac{1}{2} \left( z^n + \dfrac{1}{z^n} \right)$ and $\sin n \theta = \dfrac{1}{2i} \left(z^n - \dfrac{1}{z^n} \right)$. Thus, we have $\cos n^2 a \sin n a = \dfrac{1}{4i} \left( z^{n^2} + \dfrac{1}{z^{n^2}} \right) \left( z^{n} - \dfrac{1}{z^n} \right)$

 $= \dfrac{1}{4i} \left( z^{n^2 + n} - \dfrac{1}{z^{n^2 + n}} - z^{n^2 - n} + \dfrac{1}{z^{n^2 - n}} \right)$

 $= \dfrac{1}{2} \left( \dfrac{1}{2i} \left(z^{n^2 + n} - \dfrac{1}{z^{n^2 + n}} \right) - \dfrac{1}{2i} \left(z^{n^2 - n} - \dfrac{1}{z^{n^2 - n}} \right) \right)$

 $= \dfrac{1}{2} \left( \sin ((n^2 + n)a) - \sin ((n^2 - n)a) \right)$

 $= \dfrac{1}{2} \left( \sin(((n+1)^2 - (n+1))a) - \sin((n^2 - n)a) \right)$

 which clearly telescopes! Since the $2$ outside the brackets cancels with the $\dfrac{1}{2}$ inside, we see that the sum up to $n$ terms is 

 $\sin ((2^2 - 2)a) - \sin ((1^2 - 1)a) + \sin ((3^3 - 3)a) - \sin ((2^2 - 2)a) \cdots + \sin (((n+1)^2 - (n+1))a) - \sin ((n^2 - n)a)$

 $= \sin (((n+1)^2 - (n+1))a) - \sin(0)$

 $= \sin ((n^2 + n)a) - 0$

 $= \sin \left( \dfrac{n(n+1) \pi}{2008} \right)$.

 This expression takes on an integer value iff $\dfrac{2n(n+1)}{2008} = \dfrac{n(n+1)}{1004}$ is an integer; that is, $1004 \mid n(n+1)$. Clearly, $1004 = 2^2 \cdot 251$, implying that $251 \mid n(n+1)$. Since we want the smallest possible value of $n$, we see that we must have ${n,n+1} = 251$. If $n+1 = 251 \rightarrow n=250$, then we have $n(n+1) = 250(251)$, which is clearly not divisible by $1004$. However, if $n = 251$, then $1004 \mid n(n+1)$, so our answer is $\boxed{251}$.

 It should be noted that the product-to-sum rules follow directly from complex trigonometry, so this solution is essentially equivalent to the solution above.

 




	\item (\textit{AIME II 2008}) A particle is located on the coordinate plane at $(5,0)$. Define a \textit{move} for the particle as a counterclockwise rotation of $\pi/4$ radians about the origin followed by a translation of $10$ units in the positive $x$-direction. Given that the particle's position after $150$ moves is $(p,q)$, find the greatest integer less than or equal to $|p| + |q|$.


 
 \textbf{Answer: }019



\textbf{Solution 1}

Let $P(x, y)$ be the position of the particle on the $xy$-plane, $r$ be the length $OP$ where $O$ is the origin, and $\theta$ be the inclination of OP to the x-axis. If $(x', y')$ is the position of the particle after a move from $P$, then we have two equations for $x'$ and $y'$:
 \[x'=r\cos(\pi/4+\theta)+10 = \frac{\sqrt{2}(x - y)}{2} + 10\]
 \[y' = r\sin(\pi/4+\theta) = \frac{\sqrt{2}(x + y)}{2}.\]Let $(x_n, y_n)$ be the position of the particle after the nth move, where $x_0 = 5$ and $y_0 = 0$. Then $x_{n+1} + y_{n+1} =  \sqrt{2}x_n+10$, $x_{n+1} - y_{n+1} = -\sqrt{2}y_n+10$. This implies$x_{n+2} = -y_n + 5\sqrt{2}+ 10$, $y_{n+2}=x_n + 5\sqrt{2}$.Substituting $x_0 = 5$ and $y_0 = 0$, we have $x_8 = 5$ and $y_8 = 0$ again for the first time. Thus, $p = x_{150} = x_6 = -5\sqrt{2}$ and $q = y_{150} = y_6 = 5 + 5\sqrt{2}$. Hence, the final answer is

 $5\sqrt {2} + 5(\sqrt {2} + 1) \approx 19.1 \Longrightarrow \boxed{019}$If you're curious, the points do eventually form an octagon and repeat. Seems counterintuitive, but believe it or not, it happens.

 \href{https://www.desmos.com/calculator/febtiheosz}{https://www.desmos.com/calculator/febtiheosz}

 





\textbf{Solution 2}

Let the particle's position be represented by a complex number. Recall that multiplying a number by cis$\left( \theta \right)$ rotates the object in the complex plane by $\theta$ counterclockwise. In this case, we use $a = cis(\frac{\pi}{4})$. Therefore, applying the rotation and shifting the coordinates by 10 in the positive x direction in the complex plane results to

 $a_{150} = (((5a + 10)a + 10)a + 10 \ldots) = 5a^{150} + 10 a^{149} + 10a^{148}+ \ldots + 10$where a is cis$\left( \theta \right)$. By De-Moivre's theorem, $\left(cis( \theta \right)^n )$=cis$\left(n \theta \right)$.Therefore,

 $10(a^{150} + \ldots + 1)= 10(1 + a + \ldots + a^6) = - 10(a^7) = - 10(\frac{ \sqrt {2} }{2} - \frac{i\sqrt {2}} {2})$Furthermore, $5a^{150} = - 5i$. Thus, the final answer is

 $5\sqrt {2} + 5(\sqrt {2} + 1) \approx 19.1 \Longrightarrow \boxed{019}$





\textbf{Solution 3}

As before, consider $z$ as a complex number.  Consider the transformation $z \to (z-\omega)e^{i\theta} + \omega$.  This is a clockwise rotation of $z$ by $\theta$ radians about the points $\omega$.  Let $f(z)$ denote one move of $z$.  Then

 
                    \begin{figure}[H]
                        \centering
                        \includegraphics[width=0.5\textwidth]{Solution/AIME_2008_II_9_Solution_Figure_1.png}
                    \end{figure}
                    Therefore, $z$ rotates along a circle with center $\omega = 5+(5+5\sqrt2)i$.  Since $8 \cdot \frac{\pi}{4} = 2\pi$, $f^9(z) = f(z) \implies f^{150}(z) = f^6(z) \implies p+q = \boxed{019}$, as desired (the final algebra bash isn't bad).

 





\textbf{Solution 4}

Let $T:\begin{pmatrix}x\\y\end{pmatrix}\rightarrow R(\frac{\pi}{4})\begin{pmatrix}x\\y\end{pmatrix}+\begin{pmatrix}10\\0\end{pmatrix}$. We assume that the rotation matrix $R(\frac{\pi}{4}) = R$ here. Then we have 

 $T^{150}\begin{pmatrix}5\\0\end{pmatrix}=R(R(...R(R\begin{pmatrix}5\\0\end{pmatrix}+\begin{pmatrix}10\\0\end{pmatrix})+\begin{pmatrix}10\\0\end{pmatrix}...)+\begin{pmatrix}10\\0\end{pmatrix})+\begin{pmatrix}10\\0\end{pmatrix}$This simplifies to

 $R^{150}\begin{pmatrix}5\\0\end{pmatrix}+(I+R^2+R^3+...+R^{149})\begin{pmatrix}10\\0\end{pmatrix}$Since $R+R^{7}=O, R^2+R^6=O, R^3+R^5=O, I+R^4=O$, so we have $R^6\begin{pmatrix}5\\0\end{pmatrix}+(-R^6-R^7)\begin{pmatrix}10\\0\end{pmatrix}$, giving $p=-5\sqrt{2}, q=5\sqrt{2}+5$. The answer is yet $\lfloor10\sqrt{2}+5\rfloor=\boxed{019}$.

 




	\item (\textit{AIME II 2008}) The diagram below shows a $4\times4$ rectangular array of points, each of which is $1$ unit away from its nearest neighbors.


 
                    \begin{figure}[H]
                        \centering
                        
\includegraphics[width=0.5\textwidth]{Problem/AIME_2008_II_10_Problem_Figure_1.png}
                    \end{figure}
                    Define a \textit{growing path} to be a sequence of distinct points of the array with the property that the distance between consecutive points of the sequence is strictly increasing. Let $m$ be the maximum possible number of points in a growing path, and let $r$ be the number of growing paths consisting of exactly $m$ points. Find $mr$.


 
 \textbf{Answer: }240



\textbf{Solution}

We label our points using coordinates $0 \le x,y \le 3$, with the bottom-left point being $(0,0)$. By the \href{https://artofproblemsolving.com/wiki/index.php?title=Pythagorean\_Theorem}{Pythagorean Theorem}, the distance between two points is $\sqrt{d_x^2 + d_y^2}$ where $0 \le d_x, d_y \le 3$; these yield the possible distances (in decreasing order)
 \[\sqrt{18},\ \sqrt{13},\ \sqrt{10},\ \sqrt{9},\ \sqrt{8},\ \sqrt{5},\ \sqrt{4},\ \sqrt{2},\ \sqrt{1}\] As these define $9$ lengths, so the maximum value of $m$ is $10$. For now, we assume that $m = 10$ is achievable. Because it is difficult to immediately impose restrictions on a path with increasing distances, we consider the paths in \textit{shrinking} fashion. Note that the \textit{shrinking paths} and \textit{growing paths} are equivalent, but there are restrictions upon the locations of the first edges of the former. 

 The $\sqrt{18}$ length is only possible for one of the long diagonals, so our path must start with one of the $4$ corners of the grid. \href{https://artofproblemsolving.com/wiki/index.php?title=Without\_loss\_of\_generality}{Without loss of generality} (since the grid is rotationally symmetric), we let the vertex be $(0,0)$ and the endpoint be $(3,3)$. 

 
                    \begin{figure}[H]
                        \centering
                        \includegraphics[width=0.5\textwidth]{Solution/AIME_2008_II_10_Solution_Figure_1.png}
                    \end{figure}
                    The $\sqrt{13}$ length can now only go to $2$ points; due to reflectional symmetry about the main diagonal, we may WLOG let the next endpoint be $(1,0)$. 

 
                    \begin{figure}[H]
                        \centering
                        \includegraphics[width=0.5\textwidth]{Solution/AIME_2008_II_10_Solution_Figure_2.png}
                    \end{figure}
                    From $(1,0)$, there are two possible ways to move $\sqrt{10}$ away, either to $(0,3)$ or $(2,3)$. However, from $(0,3)$, there is no way to move $\sqrt{9}$ away, so we discard it as a possibility. 

 From $(2,3)$, the lengths of $\sqrt{8},\ \sqrt{5},\ \sqrt{4},\ \sqrt{2}$ fortunately are all determined, with the endpoint sequence being $(2,3)-(2,0)-(0,2)-(2,1)-(0,1)-(1,2)$.

 
                    \begin{figure}[H]
                        \centering
                        \includegraphics[width=0.5\textwidth]{Solution/AIME_2008_II_10_Solution_Figure_3.png}
                    \end{figure}
                    From $(1,2)$, there are $3$ possible lengths of $\sqrt{1}$ (to either $(1,1),(2,2),(1,3)$). Thus, the number of paths is $r = 4 \cdot 2 \cdot 3 = 24$, and the answer is $mr = 10 \cdot 24 = \boxed{240}$.

 




	\item (\textit{AIME II 2008}) In triangle $ABC$, $AB = AC = 100$, and $BC = 56$. Circle $P$ has radius $16$ and is tangent to $\overline{AC}$ and $\overline{BC}$. Circle $Q$ is externally tangent to circle $P$ and is tangent to $\overline{AB}$ and $\overline{BC}$. No point of circle $Q$ lies outside of $\bigtriangleup\overline{ABC}$. The radius of circle $Q$ can be expressed in the form $m - n\sqrt{k}$,where $m$, $n$, and $k$ are positive integers and $k$ is the product of distinct primes. Find $m +nk$.


 
 \textbf{Answer: }254



\textbf{Solution}


                    \begin{figure}[H]
                        \centering
                        \includegraphics[width=0.5\textwidth]{Solution/AIME_2008_II_11_Solution_Figure_1.png}
                    \end{figure}
                    Let $X$ and $Y$ be the feet of the \href{https://artofproblemsolving.com/wiki/index.php?title=Perpendicular}{perpendiculars} from $P$ and $Q$ to $BC$, respectively. Let the radius of $\odot Q$ be $r$. We know that $PQ = r + 16$. From $Q$ draw segment $\overline{QM} \parallel \overline{BC}$ such that $M$ is on $PX$. Clearly, $QM = XY$ and $PM = 16-r$. Also, we know $QPM$ is a right triangle.

 To find $XC$, consider the right triangle $PCX$. Since $\odot P$ is tangent to $\overline{AC},\overline{BC}$, then $PC$ \href{https://artofproblemsolving.com/wiki/index.php?title=Bisect}{bisects} $\angle ACB$. Let $\angle ACB = 2\theta$; then $\angle PCX = \angle QBX = \theta$. Dropping the altitude from $A$ to $BC$, we recognize the $7 - 24 - 25$ \href{https://artofproblemsolving.com/wiki/index.php?title=Right\_triangle}{right triangle}, except scaled by $4$.

 So we get that $\tan(2\theta) = 24/7$. From the \href{https://artofproblemsolving.com/wiki/index.php?title=Trigonometric\_identity}{half-angle identity}, we find that $\tan(\theta) = \frac {3}{4}$. Therefore, $XC = \frac {64}{3}$. By similar reasoning in triangle $QBY$, we see that $BY = \frac {4r}{3}$.

 We conclude that $XY = 56 - \frac {4r + 64}{3} = \frac {104 - 4r}{3}$.

 So our right triangle $QPM$ has sides $r + 16$, $r - 16$, and $\frac {104 - 4r}{3}$.

 By the \href{https://artofproblemsolving.com/wiki/index.php?title=Pythagorean\_Theorem}{Pythagorean Theorem}, simplification, and the \href{https://artofproblemsolving.com/wiki/index.php?title=Quadratic\_formula}{quadratic formula}, we can get $r = 44 - 6\sqrt {35}$, for a final answer of $\fbox{254}$.

 


 





\textbf{Solution 2}

First let $\theta = \angle{PCB}$ ; now connect the points as shown in the first solution's diagram ; realise that $\tan\theta = r/x = 16/y = r + 16/(x+y)$ where $x = BY$ and $y = CX$ (the 2 tangents) ; then we have that $QM = 64r = 56 - x - y \implies (x+y) = 56 - 64r$hence $r/x = 16+r/(56-64r)$ ; now drop altitude $AY$ to solve for $\tan{2\theta}$ ; now since we know $\tan{2\theta}$ we know $\tan \theta = r/x$ in terms of $r$ hence solve the resulting equation in $r$.

 





\textbf{Solution 3 (pure synthetic)}

Refer to the above diagram. Let the larger circle have center $O_1$, the smaller have center $O_2$, and the incenter be $I$. We can easily calculate that the area of $\triangle ABC = 2688$, and $s = 128$ and $R = 21$, where $R$ is the inradius. 

 Now, Line $\overline{AI}$ is the perpendicular bisector of $\overline{BC}$, as $\triangle ABC$ is isosceles. Letting the point of intersection be $X$, we get that $BX = 28$ and $IX = 21$, and $B, O_2, I$ are collinear as $O_2$ is equidistant from $\overline{AB}$ and $\overline{BC}$. By Pythagoras, $BI = 35$, and we notice that $\triangle BIX$ is a 3-4-5 right triangle. 

 Letting $r$ be the desired radius and letting $Y$ be the projection of $O_2$ onto $\overline{BC}$, we find that $BY = \frac{4r}{3}$. Similarly, we find that the distance between the projection from $O_1$ onto $\overline{BC}$, $W$, and $C$, is $\frac{64}{3}$. From there, we let the projection of $O_2$ onto $\overline{O_1W}$ be $Z$, and we have $O_2Z = 28 - \frac{4r}{3} + \frac{20}{3}$, $O_1Z = 16 - r$, and $O_1O_2 = 16 + r$. We finish with Pythagoras on $\triangle O_1O_2Z$, whence we get the desired answer of $\boxed{254}$. - Spacesam

 





\textbf{Solution 4}

Let the incenter be O and the altitude from A to $\overline{BC}$ be T. Note that by AA, $\triangle BQY \sim \triangle OBT$ and $\triangle PXC \sim \triangle OTC.$ Note that from $A = rs$, the inradius of the big triangle is $21$ Using ravi substitution(or noticing that $\overline{AT}$ is an altitude), we then have that $TB = TC = 28.$ From similar triangles, we can now find $\overline{BY}.$ We have 
 \[\frac{\overline{BY}}{QY} = \frac{7}{{21}} \rightarrow \overline{BY} = \frac{4}{3} r\] Now, note that as in solution 1, drawing the perpendicular from Q to $\overline{PX}$(call it Z) yields $\overline{PZ} = 16 - r, \overline{ZX} = r.$ Then, from this, 
 \[\overline{QZ} = \overline{YX} = \sqrt{(\overline{PQ})^2 - (\overline{PZ})^2} = \sqrt{(16+r)^2-(16-r)^2} = 8\sqrt{r}\] Using similar similarity as was done to find $\overline{BY}$ we have $\frac{\overline{PX}}{\overline{XC}} = \frac{\overline{OT}}{\overline{TC}} \rightarrow \frac{16}{\overline{XC}} = \frac{21}{28} \rightarrow \overline{XC} = \frac{64}{3}$.Now adding all these up and equating them to $\overline{BC}$ yields 
 \[\frac{4}{3}r + 8\sqrt{r}+ \frac{16}{3} = 56 \rightarrow r = 44 - 6\sqrt{35} \rightarrow 44 + 6\cdot 35 = \boxed{254}\]

 




	\item (\textit{AIME II 2008}) There are two distinguishable flagpoles, and there are $19$ flags, of which $10$ are identical blue flags, and $9$ are identical green flags. Let $N$ be the number of distinguishable arrangements using all of the flags in which each flagpole has at least one flag and no two green flags on either pole are adjacent. Find the remainder when $N$ is divided by $1000$.


 
 \textbf{Answer: }310



\textbf{Solution 1}

The well known problem of ordering $x$ elements of a string of $y$ elements such that none of the $x$ elements are next to each other has ${y-x+1\choose x}$ solutions. (1)

 We generalize for $a$ blues and $b$ greens. Consider a string of $a+b$ elements such that we want to choose the greens such that none of them are next to each other. We would also like to choose a place where we can divide this string into two strings such that the left one represents the first pole, and the right one represents the second pole, in order up the pole according to position on the string. 

 However, this method does not account for the fact that the first pole could end with a green, and the second pole could start with a green, since the original string assumed that no greens could be consecutive. We solve this problem by introducing an extra blue, such that we choose our divider by choosing one of these $a+1$ blues, and not including that one as a flag on either pole.

 From (1), we now have ${a+2\choose b}$ ways to order the string such that no greens are next to each other, and $a+1$ ways to choose the extra blue that will divide the string into the two poles: or $(a+1){a+2\choose b}$ orderings in total.

 However, we have overcounted the solutions where either pole has no flags, we have to count these separately. This is the same as choosing our extra blue as one of the two ends, and ordering the other $a$ blues and $b$ greens such that no greens are next to each other: for a total of $2{a+1\choose b}$ such orderings.

 Thus, we have $(a+1){a+2\choose b}-2{a+1\choose b}$ orderings that satisfy the conditions in the problem: plugging in $a=10$ and $b=9$, we get $2310 \equiv \boxed{310} \pmod{1000}$.

 





\textbf{Solution 2}

Split the problem into two cases:

 Case 1 - Both poles have blue flags:There are 9 ways to place the 10 blue flags on the poles. In each of these configurations, there are 12 spots that a green flag could go. (either in between two blues or at the tops or bottoms of the poles) Then, since there are 9 green flags, all of which must be used, there are ${12\choose9}$ possiblities for each of the 9 ways to place the blue flags. Total: ${12\choose9}*9$ possibilities.

 Case 2 - One pole has no blue flags:Since each pole is non empty, the pole without a blue flag must have one green flag. The other pole has 10 blue flags and, by the argument used in case 1, there are ${11\choose8}$ possibilities, and since the poles are distinguishable, there are a total of $2*{11\choose8}$ possiblities for this case.

 Finally, we have $9*{12\choose9}+2*{11\choose8}=2310 \equiv \boxed{310} \pmod{1000}$ as our answer.

 





\textbf{Solution 3 (mad bash)}

Call the two flagpoles Flagpole A and Flagpole B.Case 1:Flag distribution: 1|18B|variable: $\dbinom{10}{9}=10$ waysG|variable: $\dbinom{11}{8}=165$ waysCase 2:Flag distribution: 2|17BB|variable: $\dbinom{9}{9}=1$ wayBG|variable (two ways to arrange the first sequence): $\dbinom{10}{8} \times 2 = 90$ waysGG|variable: can't happenCase 3: Flag distribution: 3|16BBB|variable: can't happenBBG|variable (3 legal ways to arrange): $\dbinom{9}{8} \times 3 = 27$ waysGBG|variable (1 legal way to arrange): $\dbinom{10}{7} = 120$ waysGGG|variable: clearly can't happenCase 4:Flag distribution: 4|15BBBB|variable: can't happenBBBG|variable (4 legal ways to arrange): $\dbinom{8}{8} \times 4 = 4$ waysGBBG|variable (3 legal ways to arrange): $\dbinom{9}{7} \times 3 = 108$ waysGGBG|variable: can't happenGGGG|variable: can't happenCase 5:Flag distribution: 5|14BBBBB|variable: can't happenBBBBG|variable: can't happenBBGBG|variable (6 legal ways to arrange): $\dbinom{8}{7} \times 6 = 48$ waysGBGBG|variable (1 legal way to arrange): $\dbinom{9}{6} = 84$ waysGGGGB|variable: can't happenGGGGG|variable: can't happen Case 6:Flag distribution: 6|13BBBBBB|variable: can't happenBBBBBG|variable: can't happenBBBBGG|variable (10 legal ways to arrange): $\dbinom{7}{7} \times 10 = 10$ waysBBBGGG|variable (4 legal ways to arrange): $\dbinom{8}{6} \times 4 = 112$ waysBBGGGG|variable: can't happenBGGGGG|variable: can't happenGGGGGG|variable: can't happenCase 7: Flag distribution: 7|12BBBBBBB|variable: can't happenBBBBBBG|variable: can't happenBBBBBGG|variable: can't happenBBBBGGG|variable (10 legal ways to arrange): $\dbinom{7}{6} \times 10 = 70$ waysBBBGGGG|variable (1 legal way to arrange): $\dbinom{8}{5} = 56$ waysrest can't happenCase 8:Flag distribution: 8|11BBBBBBBB|variable: can't happenBBBBBBBG|variable: can't happenBBBBBBGG|variable: can't happenBBBBBGGG|variable (20 legal ways to arrange): $\dbinom{6}{6} \times 20 = 20$ waysBBBBGGGG|variable: (5 legal ways to arrange): $\dbinom{7}{5} \times 5 = 105$ waysothers can't happenCase 9:Flag distribution: 9|10BBBBBGGGG|variable (15 legal ways to arrange): $\dbinom{6}{5} \times 15 = 90$ waysBBBBGGGGG|variable (1 legal way to arrange): $\dbinom{7}{4} = 35$ waysothers can't happenSo our total number of ways is $2$ times $10+165+1+90+27+120+4+108+48+84+10+112+70+56+20+105+90+35$ (since the two flagpoles are distinguishable) which is $2310$ ways. We have to find the last 3 digits, so our final answer is $310$.

 Solution by fidgetboss\_4000

 Note: Do not attempt unless you are good at bashing.

 


 


 





\textbf{Solution 4 (Vandermonde's)}

Let the number of blue flags on the first flagpole be $b$ and define $g$ similarly. Now note the bound $g \le b+1.$ Now we cannot have any two consecutive greens. This condition is the same as "there are $b+1$ spaces between each blue flag. How many ways can we put the $g$ green flags in the $b+1$ spaces?" Therefore given $b$ blue flags on the first flagpole and $g$ green flags on the first flagpole we have a total of $\binom{b+1}{g}.$ Similarly there are $\binom{11-b}{9-g}$ ways for the second flagpole. Summing over all possible possibilities we see the sum is $\sum_{g = 0}^{b+1} \sum_{b=0}^{10}\binom{b+1}{g}\binom{11-b}{9-g}.$ Swapping the sum we have $\sum_{g = 0}^{b+1} \sum_{b=0}^{10}\binom{b+1}{g}\binom{11-b}{9-g} = \sum_{b=0}^{10} \sum_{g = 0}^{b+1}\binom{b+1}{g}\binom{11-b}{9-g}.$ Then applying Vandermonde's yields $\sum_{b=0}^{10}\binom{12}{9}$ so the sum is $11 \cdot\binom{12}{9}.$ However, this answer overcounts. We cannot have no flags on a pole, so we subtract $2 \cdot\binom{11}{9}$ since we can have empty flags for the first or second flag. Therefore the answer is $11 \cdot\binom{12}{9} - 2 \cdot\binom{11}{9} \pmod{1000} \equiv \boxed{310} \pmod{1000}.$

 





\textbf{Solution 5}

Consider this as arranging the flags on one big flagpole, then splitting the flagpole into two. 

 Example: Start with big flagpole $BGBGBGGBBGBGBGBGBGB$, and then split it into $BGBGBG$ and $GBBGBGBGBGBGB$.

 We will split this problem into two cases:

 Case 1: The split is not two greens. Basically if arranging 10 blue flags and 9 green flags does not give any consecutive green flags, then we can split it into two flagpoles wherever. For example,
 \[BGBBGBGBGBGBGBGBGBG\]can become 
 \[BGB|BGBGBGBGBGBGBGBG\] or 
 \[BGBBG|BGBGBGBGBGBGBG\].

 By stars and bars, there are $\binom{11}{2}$ ways to arrange the flags, then $18$ ways to split the big flagpole, so total $\binom{11}{2}\cdot 18$ ways.

 Case 2: The split is two greens. Here the big flagpole has two consecutive green flags, so we are forced to split the two green flags.

 There are $\binom{11}{3}\cdot 8$ ways for this case, since we first place 8 “$G$”s among 10 blue flags, then choose one to become $GG$. (Look at the first example)

 Thus the total number of ways is
 \[\binom{11}{2}\cdot 18 + \binom{11}{3}\cdot 8 = 2310 \equiv \boxed{310} \pmod{1000}.\]

 





\textbf{Solution 6}

We split the 9 green flags over the two poles: $0$ and $9$, $1$ and $8$, $2$ and $7$, and so on, respectively.

 For the first case, $0$ and $9$, 8 of the 10 blue flags must be placed between the 9 green flags on the same pole, and the 9th blue flag must be placed on the other pole so that there is at least one flag on each pole. This leaves $\textbf{11}$ choices for the last blue flag.

 Moving on, notice that now for every other case there will always be 3 blue flags left to place. Since we want to distribute 3 indistinguishable flags to 11 distinguishable spots, we use stars and bars. There are $\binom{13}{10} = \textbf{286}$ ways to distribute the flags for the other cases.

 Then the total number of ways is $2 * (11 + 4 * 286) = 2310 \equiv \boxed{310} \pmod{1000}$, since the poles are distinguishable.

 





\textbf{Solution 7 (Fast and relatively simple)}

Say we combine the two flagpoles such that we have a line with 10 blue flags and 9 green flags. To distinguish the two flagpoles from each other, we add a "separator" that acts as both a way to distinguish the two flagpoles and a way to show that the separator can split green flags too. Therefore, we now have 10 blue flags, 9 green flags, and 1 separator. We first place the 10 blue flags and 1 separator. There are $\binom{11}{1}$ ways to accomplish this, and these flags establish the walls between green flags. Next, we want to place the green flags such that they aren't adjacent. We must place the greens in between other flags, and because we have 11 non-green flags already placed, we have $\binom{12}{9}$ ways to place the green flags. Now, we must subtract the cases where a flagpole has no flags. This happens when we place the separator on the ends, as in this case: SGBGBGBGBGBGBGBGBGBB. There are $2$ ends, and $\binom{11}{9}$ ways to arrange the remaining 9 greens and 10 blues. Thus, our final answer is $\binom{11}{1} * \binom{12}{9} - 2 * \binom{11}{9}$ = $2310$. Taking mod 1000, we get that the final answer is $\boxed{310}$.

 




	\item (\textit{AIME II 2008}) A regular hexagon with center at the origin in the complex plane has opposite pairs of sides one unit apart. One pair of sides is parallel to the imaginary axis. Let $R$ be the region outside the hexagon, and let $S = \left\lbrace\frac{1}{z}|z \in R\right\rbrace$. Then the area of $S$ has the form $a\pi + \sqrt{b}$, where $a$ and $b$ are positive integers. Find $a + b$.


 
 \textbf{Answer: }029



\textbf{Solution 1}

If you're familiar with inversion, you'll see that the problem's basically asking you to invert the hexagon with respect to the unit circle in the Cartesian Plane using the Inversion Distance Formula. This works because the point in the Cartesian Plane's complex plane equivalent switches places with its conjugate but we can do that in the Cartesian plane too (just reflect a point in the Cartesian plane over the x-axis)! If you're familiar with inversion you can go plot the inverted figure's Cartesian Plane Equivalent. Then simply continue on with the figure shown in the below solution.

 





\textbf{Solution 2}

If a point $z = r\text{cis}\,\theta$ is in $R$, then the point $\frac{1}{z} = \frac{1}{r} \text{cis}\, \left(-\theta\right)$ is in $S$ (where \href{https://artofproblemsolving.com/wiki/index.php?title=Cis}{cis} denotes $\text{cis}\, \theta = \cos \theta + i \sin \theta$). Since $R$ is symmetric every $60^{\circ}$ about the origin, it suffices to consider the area of the result of the transformation when $-30 \le \theta \le 30$, and then to multiply by $6$ to account for the entire area.

 We note that if the region $S_2 = \left\lbrace\frac{1}{z}|z \in R_2\right\rbrace$, where $R_2$ is the region (in green below) outside the circle of radius $1/\sqrt{3}$ centered at the origin, then $S_2$ is simply the region inside a circle of radius $\sqrt{3}$ centered at the origin. It now suffices to find what happens to the mapping of the region $R-R_2$ (in blue below). 

 The equation of the hexagon side in that region is $x = r \cos \theta = \frac{1}{2}$, which is transformed to $\frac{1}{r} \cos (-\theta) = \frac{1}{r} \cos \theta =$2 . Let $r\text{cis}\,\theta = a+bi$ where $a,b \in \mathbb{R}$; then $r = \sqrt{a^2 + b^2}, \cos \theta = \frac{a}{\sqrt{a^2 + b^2}}$, so the equation becomes $a^2 - 2a + b^2 = 0 \Longrightarrow (a-1)^2 + b^2 = 1$. Hence the side is sent to an arc of the unit circle centered at $(1,0)$, after considering the restriction that the side of the hexagon is a segment of length $1/\sqrt{3}$. 

 Including $S_2$, we find that $S$ is the union of six unit circles centered at $\text{cis}\, \frac{k\pi}{6}$, $k = 0,1,2,3,4,5$, as shown below.  

 
                    \begin{figure}[H]
                        \centering
                        \includegraphics[width=0.5\textwidth]{Solution/AIME_2008_II_13_Solution_Figure_1.png}
                    \end{figure}
                         $\Longrightarrow$     
                    \begin{figure}[H]
                        \centering
                        \includegraphics[width=0.5\textwidth]{Solution/AIME_2008_II_13_Solution_Figure_2.png}
                    \end{figure}
                    The area of the regular hexagon is $6 \cdot \left( \frac{\left(\sqrt{3}\right)^2 \sqrt{3}}{4} \right) = \frac{9}{2}\sqrt{3}$. The total area of the six $120^{\circ}$ sectors is $6\left(\frac{1}{3}\pi - \frac{1}{2} \cdot \frac{1}{2} \cdot \sqrt{3}\right) = 2\pi - \frac{3}{2}\sqrt{3}$. Their sum is $2\pi + \sqrt{27}$, and $a+b = \boxed{029}$.- Th3Numb3rThr33

 





\textbf{Solution 3 (Calculus)}

We can describe the line parallel to the imaginary axis $x=\frac{1}{2}$ using polar coordinates as $r(\theta)=\dfrac{1}{2\cos{\theta}},$

 which rearranges to $z=\left(\dfrac{1}{2\cos{\theta}}\right)(cis{\theta})\implies \frac{1}{z}=2\cos{\theta}cis(-\theta).$

 Denote the area of $S$ as $[S].$ Now, dividing the hexagon to 12 equal parts, we have the following integral:

 $[S] = 12\int_{0}^{\frac{\pi}{6}}\frac{1}{2}r^2 d\theta = 12\int_{0}^{\frac{\pi}{6}}\frac{1}{2}(2\cos\theta)^2 d\theta.$

 Thankfully, this is a routine computation:

 $[S] = 12\int_{0}^{\frac{\pi}{6}}2(\cos\theta)^2 d\theta = 12\int_{0}^{\frac{\pi}{6}}(\cos{2\theta}+1)d\theta$

 $[S] = 12\int_{0}^{\frac{\pi}{6}}(\cos{2\theta}+1)d\theta = 12\left[\frac{1}{2}\sin{2\theta}+\theta\right]_0^{\frac{\pi}{6}}=12\left(\frac{\sqrt{3}}{4}+\frac{\pi}{6}\right)=2\pi+3\sqrt{3}=2\pi + \sqrt{27}$

 $a+b = \boxed{029}$.

 


 \href{https://mathproblemsolvingskills.wordpress.com/}{MathProblemSolvingSkills.com}

 


 




	\item (\textit{AIME II 2008}) Let $a$ and $b$ be positive real numbers with $a \ge b$. Let $\rho$ be the maximum possible value of $\dfrac{a}{b}$ for which the system of equations

 \[a^2 + y^2 = b^2 + x^2 = (a-x)^2 + (b-y)^2\]
has a solution $(x,y)$ satisfying $0 \le x < a$ and $0 \le y < b$. Then $\rho^2$ can be expressed as a fraction $\dfrac{m}{n}$, where $m$ and $n$ are relatively prime positive integers. Find $m + n$.


 
 \textbf{Answer: }007



\textbf{Solution 1}

Notice that the given equation implies

 $a^2 + y^2 = b^2 + x^2 = 2(ax + by)$We have $2by \ge y^2$, so $2ax \le a^2 \implies x \le \frac {a}{2}$.

 Then, notice $b^2 + x^2 = a^2 + y^2 \ge a^2$, so $b^2 \ge \frac {3}{4}a^2 \implies \rho^2 \le \frac {4}{3}$.

 The solution $(a, b, x, y) = \left(1, \frac {\sqrt {3}}{2}, \frac {1}{2}, 0\right)$ satisfies the equation, so $\rho^2 = \frac {4}{3}$, and the answer is $3 + 4 = \boxed{007}$.

 





\textbf{Solution 2}

Consider the points $(a,y)$ and $(x,b)$. They form an \href{https://artofproblemsolving.com/wiki/index.php?title=Equilateral\_triangle}{equilateral triangle} with the origin. We let the side length be $1$, so $a = \cos{\theta}$ and $b = \sin{\left(\theta + \frac {\pi}{3}\right)}$.

 Thus $f(\theta) = \frac {a}{b} = \frac {\cos{\theta}}{\sin{\left(\theta + \frac {\pi}{3}\right)}}$ and we need to maximize this for $0 \le \theta \le \frac {\pi}{6}$.

 Taking the \href{https://artofproblemsolving.com/wiki/index.php?title=Derivative}{derivative} shows that $-f'(\theta) = \frac {\cos{\frac {\pi}{3}}}{\sin^2{\left(\theta + \frac {\pi}{3}\right)}} \ge 0$, so the maximum is at the endpoint $\theta = 0$. We then get

 $\rho = \frac {\cos{0}}{\sin{\frac {\pi}{3}}} = \frac {2}{\sqrt {3}}$Then, $\rho^2 = \frac {4}{3}$, and the answer is $3+4=\boxed{007}$.

 (For a non-calculus way to maximize the function above:

 Let us work with degrees. Let $f(x)=\frac{\cos x}{\sin(x+60)}$. We need to maximize $f$ on $[0,30]$. 

 Suppose $k$ is an upper bound of $f$ on this range; in other words, assume $f(x)\le k$ for all $x$ in this range. Then: 
 \[\cos x\le k\sin(x+60)=k\cdot\left(\frac{\sqrt{3}}{2}\cos x+\frac{1}{2}\sin x\right)\]
 \[\rightarrow 0\le \left(\frac{\sqrt{3}k}{2}-1\right)\cos x+\frac{k}{2}\sin x\rightarrow 0\le (\sqrt{3}k-2)\cos x+k\sin x\]
 \[\rightarrow (2-\sqrt{3}k)\cos x\le k\sin x\rightarrow \frac{2-\sqrt{3}k}{k}\le \tan x,\]for all $x$ in $[0,30]$. In particular, for $x=0$, $\frac{2-\sqrt{3}k}{k}$ must be less than or equal to $0$, so $k\ge \frac{2}{\sqrt{3}}$.

 The least possible upper bound of $f$ on this interval is $k=\frac{2}{\sqrt{3}}$. This inequality must hold by the above logic, and in fact, the inequality reaches equality when $x=0$. Thus, $f(x)$ attains a maximum of $\frac{2}{\sqrt{3}}$ on the interval.)

 





\textbf{Solution 3}

Consider a \href{https://artofproblemsolving.com/wiki/index.php?title=Cyclic\_quadrilateral}{cyclic quadrilateral} $ABCD$ with $\angle B = \angle D = 90^{\circ}$, and $AB = y, BC = a, CD = b, AD = x$. Then
 \[AC^2 = a^2 + y^2 = b^2 + x^2\]From \href{https://artofproblemsolving.com/wiki/index.php?title=Ptolemy\%27s\_Theorem}{Ptolemy's Theorem}, $ax + by = AC(BD)$, so
 \[AC^2 = (a - x)^2 + (b - y)^2 = a^2 + y^2 + b^2 + x^2 - 2(ax + by) = 2AC^2 - 2AC*BD\]Simplifying, we have $BD = AC/2$. 

 Note the \href{https://artofproblemsolving.com/wiki/index.php?title=Circumcircle}{circumcircle} of $ABCD$ has \href{https://artofproblemsolving.com/wiki/index.php?title=Radius}{radius} $r = AC/2$, so $BD = r$ and has an arc of $60^{\circ}$, so $\angle C = 30^{\circ}$. Let $\angle BDC = \theta$.

 $\frac ab = \frac{BC}{CD} = \frac{\sin \theta}{\sin(150^{\circ} - \theta)}$, where both $\theta$ and $150^{\circ} - \theta$ are $\leq 90^{\circ}$ since triangle $BCD$ must be \href{https://artofproblemsolving.com/wiki/index.php?title=Acute\_triangle}{acute}. Since $\sin$ is an increasing function over $(0, 90^{\circ})$, $\frac{\sin \theta}{\sin(150^{\circ} - \theta)}$ is also increasing function over $(60^{\circ}, 90^{\circ})$.

 $\frac ab$ maximizes at $\theta = 90^{\circ} \Longrightarrow \frac ab$ maximizes at $\frac 2{\sqrt {3}}$. This squared is $(\frac 2{\sqrt {3}})^2 = \frac4{3}$, and $4 +  3 = \boxed{007}$.

 \textbf{Note:} None of the above solutions point out clearly the importance of the restriction that $a$, $b$, $x$ and $y$ be positive. Indeed, larger values of p are obtained when the lower vertex of the equilateral triangle in Solution 2 dips below the x-axis. Take for example $-15= \theta$. This yields $p = (1 + \sqrt{3})/2 > 4/3$

 





\textbf{Solution 4}

The problem is looking for an intersection in the said range between parabola $P$: $y = \frac{(x-a)^2 + b^2-a^2}{2b}$ and the hyperbola $H$: $y^2 = x^2 + b^2 - a^2$. The vertex of $P$ is below the x-axis and it's x-coordinate is a, which is to the right of the vertex of the $H$, which is $\sqrt{a^2 - b^2}$. So for the intersection to exist with $x<a$ and $y \geq 0$, $P$ needs to cross x-axis between $\sqrt{a^2 - b^2}$, and $a$, meaning,
 \[(\sqrt{a^2 - b^2}-a)^2 + b^2-a^2 \geq 0\]Divide both side by $b^2$,
 \[(\sqrt{\rho^2 - 1}-\rho)^2 + 1-\rho^2 \geq 0\]which can be easily solved by moving $1-\rho^2$ to RHS and taking square roots. Final answer $\rho^2 \leq \frac{4}{3}$$\boxed{007}$

 





\textbf{Solution 5}

The given system is equivalent to the points $(a,y)$ and $(x,b)$ forming an equilateral triangle with the origin. WLOG let this triangle have side length $1$, so $x=\sqrt{1-a^2}$. The condition $x<a$ implies $(x,b)$ lies to the left of $(a,y)$, so $(x,b)$ is the top vertex. Now we can compute (by complex numbers, or the sine angle addition identity) that $b = \frac{\sqrt{3}}{2}a + \frac{1}{2}\sqrt{1-a^2}$, so $\frac{a}{b} = \frac{a}{\frac{\sqrt{3}}{2}a + \frac{1}{2}\sqrt{1-a^2}} = \frac{1}{\frac{\sqrt{3}}{2} + \frac{1}{2a}\sqrt{1-a^2}}$. Minimizing this is equivalent to minimizing the denominator, which happens when $\sqrt{1-a^2} = 0$ and thus $a=1$, resulting in $\rho = \frac{2}{\sqrt{3}}$, so $\rho^2 = \frac{4}{3}$ and the answer is $\boxed{007}$.

 As a remark, expressing the condition that the triangle is equilateral purely algebraically instead of using trig eliminates the need for calculus or analyzing the behavior of sine.

 





\textbf{Solution 6 (Geometry and Trigonometry)}

Notice that by Pythagorean theorem, if we take a triangle with vertices $(0,0),$ $(a,y),$ and $(x,b)$ forming an equilateral triangle. Now, take a rectangle with vertices $(0,0), (a,0), (0,b), (a,b).$ Notice that $(a,y)$ and $(x,b)$ are on the sides. Let $\alpha$ be the angle formed by the points $(0,b), (0,0), (x,b).$ Then, we have that 
 \[\cos \alpha = \frac{b}{s},\] where $s$ is the side of the equilateral triangle. Also, we have that $30^{\circ}-\alpha$ is the angle formed by the points $(a,0), (0,0), (a,y),$ and so 
 \[\cos (30^{\circ}-\alpha) = \frac{a}{s}.\] Thus, we have that 
 \[\frac{a}{b} = \frac{\cos (30^{\circ}-\alpha)}{\cos \alpha}.\] We see that this expression is maximized when $\alpha$ is maximized (at least when $\alpha$ is in the interval $(0,90^{\circ}),$ which it is). Then, $\alpha \ge 30^{\circ},$ so ew have that the maximum of $\frac{a}{b}$ is 
 \[\frac{\cos 0}{\cos 30^{\circ}} = \frac{2}{\sqrt{3}},\] and so our answer is $4+3 = 7.$

 





\textbf{Solution 7 (No calc/trig)}

As the previous solutions pointed out, we have that the triangle with vertices $(0,0),$ $(a,y),$ and $(x,b)$ is equilateral. Then, note that that equilateral triangle is inscribed in a rectangle with vertices $(0,0),$ $(a,0),$ $(a,b),$ and $(0,b).$ Thus, our objective is to find the most "oblique" rectangle that has an equilateral triangle inscribed in it. 

 This is when a side of the equilateral triangle coincides with a side of the rectangle, and therefore an axis. We see this because if we rotate it so that it is no longer the case, the longer side decreases in length while the shorter side increases in length. Rotating it to coincides with the other axis, it is just the original rectangle reflected over $y=x.$ WLOG let the side length of the equilateral triangle be $1.$ Then, by our above argument, 
 \[(a,b)=\left(1,\dfrac{\sqrt{3}}{2}\right)\] is optimal, so our answer is
 \[\left(\dfrac{2}{\sqrt{3}}\right)^2=\dfrac{4}{3}\implies \boxed{007}.\]

 


 \href{https://mathproblemsolvingskills.wordpress.com/}{MathProblemSolvingSkills.com}

 


 




	\item (\textit{AIME II 2008}) Find the largest integer $n$ satisfying the following conditions:


 
            \begin{enumerate}[label=(\roman*)]
                \item  $n^2$ can be expressed as the difference of two consecutive cubes;
\item  $2n + 79$ is a perfect square.
            \end{enumerate}
            
 \textbf{Answer: }181



\textbf{Solution 1}

Write $n^2 = (m + 1)^3 - m^3 = 3m^2 + 3m + 1$, or equivalently, $(2n + 1)(2n - 1) = 4n^2 - 1 = 12m^2 + 12m + 3 = 3(2m + 1)^2$.

 Since $2n + 1$ and $2n - 1$ are both odd and their difference is $2$, they are \href{https://artofproblemsolving.com/wiki/index.php?title=Relatively\_prime}{relatively prime}. But since their product is three times a square, one of them must be a square and the other three times a square. We cannot have $2n - 1$ be three times a square, for then $2n + 1$ would be a square congruent to $2$ modulo $3$, which is impossible.

 Thus $2n - 1$ is a square, say $b^2$. But $2n + 79$ is also a square, say $a^2$. Then $(a + b)(a - b) = a^2 - b^2 = 80$. Since $a + b$ and $a - b$ have the same parity and their product is even, they are both even. To maximize $n$, it suffices to maximize $2b = (a + b) - (a - b)$ and check that this yields an integral value for $m$. This occurs when $a + b = 40$ and $a - b = 2$, that is, when $a = 21$ and $b = 19$. This yields $n = 181$ and $m = 104$, so the answer is $\boxed{181}$.

 





\textbf{Solution 2}

Suppose that the consecutive cubes are $m$ and $m + 1$. We can use completing the square and the first condition to get:
 \[(2n)^2 - 3(2m + 1)^2 = 1\equiv a^2 - 3b^2\]where $a$ and $b$ are non-negative integers. Now this is a \href{https://artofproblemsolving.com/wiki/index.php?title=Pell\_equation}{Pell equation}, with solutions in the form $(2 + \sqrt {3})^k = a_k + \sqrt {3}b_k,$ $k = 0,1,2,3,...$. However, $a$ is even and $b$ is odd. It is easy to see that the parity of $a$ and $b$ switch each time (by induction). Hence all solutions to the first condition are in the form:
 \[(2 + \sqrt {3})^{2k + 1} = a_k + \sqrt {3}b_k\]where $k = 0,1,2,..$. So we can (with very little effort) obtain the following: $(k,a_k = 2n) = (0,2),(1,26),(2,362),(3,5042)$. It is an AIME problem so it is implicit that $n < 1000$, so $2n < 2000$. It is easy to see that $a_n$ is strictly increasing by induction. Checking $2n = 362\implies n =\boxed{181}$ in the second condition works (we know $b_k$ is odd so we don't need to find $m$). So we're done.

 





\textbf{Solution 3 (Big Bash)}

Let us generate numbers $1$ to $1000$ for the second condition, for squares. We know for $N$ to be integer, the squares must be odd. So we generate $N = 1, 21, 45, 73, 105, 141, 181, 381, 441, 721, 801$. $N$ cannot exceed $1000$ since it is AIME problem. Now take the first criterion, let $a$ be the smaller consecutive cube. We then get:

 $N^2 = (A + 1)^3 - A^3$

 $N^2 - 1 = 3A^2 + 3A$

 $(N + 1)(N - 1) = 3A(A + 1)$

 Now we know either $N + 1$ or $N - 1$ must be factor of $3$, hence $N = 1 \pmod 3$ or$N = 2 \pmod 3$. Only $1, 73, 181, 721$ satisfy this criterion. Testing each of the numbers in the condition yields $181$ as the largest that fits both, thus answer $= \boxed{181}$.

 


 \href{https://mathproblemsolvingskills.wordpress.com/}{MathProblemSolvingSkills.com}

 


 




\end{enumerate}